\documentclass[10pt,a4paper]{amsart}
\usepackage{amsmath,amssymb,amsthm,graphicx}
% Retained preamble for the Stacks Project chapter snapshot in this directory.
% The chapter's own preamble.tex is a copy of the shared Stacks Project preamble; the block
% below is sliced MECHANICALLY, line for line, from it. Nothing was added, removed or
% redefined; the wrapper class, the xr-hyper/hyperref wiring and the externaldocument list
% are omitted because the assembler supplies the document class.
%
% lines taken from preamble.tex: 3, 5, 9, 11, 12, 156, 157, 158, ... (51 lines total)
%
% package coverage restored from the chapter's own preamble (2026-10-13): xy, xy-2cell

\usepackage[all]{xy}
\xyoption{2cell}
\UseAllTwocells

\newenvironment{reference}{\comment}{\endcomment}
\newenvironment{slogan}{\comment}{\endcomment}
\newenvironment{history}{\comment}{\endcomment}
\theoremstyle{plain}
\newtheorem{theorem}[subsection]{Theorem}
\newtheorem{proposition}[subsection]{Proposition}
\newtheorem{lemma}[subsection]{Lemma}
\theoremstyle{definition}
\newtheorem{definition}[subsection]{Definition}
\newtheorem{example}[subsection]{Example}
\newtheorem{exercise}[subsection]{Exercise}
\newtheorem{situation}[subsection]{Situation}
\theoremstyle{remark}
\newtheorem{remark}[subsection]{Remark}
\newtheorem{remarks}[subsection]{Remarks}
\numberwithin{equation}{subsection}
\def\lim{\mathop{\mathrm{lim}}\nolimits}
\def\colim{\mathop{\mathrm{colim}}\nolimits}
\def\Spec{\mathop{\mathrm{Spec}}}
\def\Hom{\mathop{\mathrm{Hom}}\nolimits}
\def\Ext{\mathop{\mathrm{Ext}}\nolimits}
\def\SheafHom{\mathop{\mathcal{H}\!\mathit{om}}\nolimits}
\def\SheafExt{\mathop{\mathcal{E}\!\mathit{xt}}\nolimits}
\def\Sch{\mathit{Sch}}
\def\Mor{\mathop{\mathrm{Mor}}\nolimits}
\def\Ob{\mathop{\mathrm{Ob}}\nolimits}
\def\Sh{\mathop{\mathit{Sh}}\nolimits}
\def\NL{\mathop{N\!L}\nolimits}
\def\CH{\mathop{\mathrm{CH}}\nolimits}
\def\proetale{{pro\text{-}\acute{e}tale}}
\def\etale{{\acute{e}tale}}
\def\QCoh{\mathit{QCoh}}
\def\Ker{\mathop{\mathrm{Ker}}}
\def\Im{\mathop{\mathrm{Im}}}
\def\Coker{\mathop{\mathrm{Coker}}}
\def\Coim{\mathop{\mathrm{Coim}}}
\DeclareMathSymbol{\boxtimes}{\mathbin}{AMSa}{"02}
\def\QCohstack{\mathcal{QC}\!\mathit{oh}}
\def\Cohstack{\mathcal{C}\!\mathit{oh}}
\def\Spacesstack{\mathcal{S}\!\mathit{paces}}
\def\Quotfunctor{\mathrm{Quot}}
\def\Hilbfunctor{\mathrm{Hilb}}
\def\Curvesstack{\mathcal{C}\!\mathit{urves}}
\def\Polarizedstack{\mathcal{P}\!\mathit{olarized}}
\def\Complexesstack{\mathcal{C}\!\mathit{omplexes}}
\def\Pic{\mathop{\mathrm{Pic}}\nolimits}
\def\Picardstack{\mathcal{P}\!\mathit{ic}}
\def\Picardfunctor{\mathrm{Pic}}
\def\Deformationcategory{\mathcal{D}\!\mathit{ef}}
\usepackage{cleveref}
\crefname{originallemma}{original lemma}{original lemmas}
\Crefname{originallemma}{Original lemma}{Original lemmas}
\makeatletter
\newlabel{simplicial-lemma-product-trivial-kan}{{08NR}{source}{}{sourcepointer-1}{}}
\newlabel{simplicial-lemma-product-trivial-kan@cref}{{[originallemma][0][]08NR}{[1][0][]source}{}{}{}}
\makeatother
\crefname{originallemma}{original lemma}{original lemmas}
\Crefname{originallemma}{Original lemma}{Original lemmas}
\makeatletter
\newlabel{simplicial-lemma-trivial-kan-base-change}{{08NN}{source}{}{sourcepointer-2}{}}
\newlabel{simplicial-lemma-trivial-kan-base-change@cref}{{[originallemma][0][]08NN}{[1][0][]source}{}{}{}}
\makeatother
\crefname{originallemma}{original lemma}{original lemmas}
\Crefname{originallemma}{Original lemma}{Original lemmas}
\makeatletter
\newlabel{simplicial-lemma-trivial-kan-homotopy}{{08NS}{source}{}{sourcepointer-3}{}}
\newlabel{simplicial-lemma-trivial-kan-homotopy@cref}{{[originallemma][0][]08NS}{[1][0][]source}{}{}{}}
\makeatother
\crefname{originallemma}{original lemma}{original lemmas}
\Crefname{originallemma}{Original lemma}{Original lemmas}
\makeatletter
\newlabel{sites-cohomology-lemma-compute-by-cosimplicial-resolution}{{08Q9}{source}{}{sourcepointer-4}{}}
\newlabel{sites-cohomology-lemma-compute-by-cosimplicial-resolution@cref}{{[originallemma][0][]08Q9}{[1][0][]source}{}{}{}}
\makeatother
\title{The lower shriek computed by a polynomial simplicial resolution with trivial Kan fibration augmentation}
\author{The Stacks Project Authors}
\date{stacks-project:cotangent.tex:local-snapshot-20261012; GNU FDL 1.2 or later}
\begin{document}
\maketitle
\setcounter{subsection}{1}
\noindent\textit{Editorial note.} All delivered mathematical text is the author's own, copied from the pinned Stacks Project chapter file cotangent.tex; no statement or proof step was rewritten, reordered or removed.\par\medskip
\noindent\textit{Editorial note.} The statement carries two displays (the isomorphism and the reference to the functoriality in F); the proof works inside the category C fixed earlier in the chapter.\par\medskip
\noindent\textit{Editorial note.} The four references of the counted spans are recorded as non-load-bearing navigation pointers; the argument that consumes them stays inside the counted proof span.\par\medskip
\noindent\textit{Editorial note.} Preamble: the retained preamble of this chapter (sources/stacks-cotangent/retained-preamble.tex, lines 1-63), containing the Xy-pic setup, the comment environments, the theorem declarations and the macros including \textbackslash{}NL.\par\medskip
\noindent\textit{Allowed standard prerequisites (editorial).}\begin{itemize}
\item The cosimplicial resolution criterion for computing lower shriek. Cohomology on Sites, Lemma sites-cohomology-lemma-compute-by-cosimplicial-resolution, used at proof lines 458-465. Reduces the claim to the homotopy statement about S\_bullet.
\item Stability of trivial Kan fibrations under base change and finite products. Simplicial, Lemmas simplicial-lemma-trivial-kan-base-change and simplicial-lemma-product-trivial-kan, used at proof lines 481-485. Give that the fibres and their product over E are trivial Kan fibrations.
\item A trivial Kan fibration is homotopy equivalent to a point. Simplicial, Lemma simplicial-lemma-trivial-kan-homotopy, used at proof lines 486-487. Concludes the homotopy statement.
\item Every object of the category C is a polynomial algebra over A. used at proof lines 466-467. Allows writing Q = A[E].
\end{itemize}\medskip
\ifdefined\hypertarget\hypertarget{sourcepointer-1}{}\fi
\noindent\textit{Source reference.} Original lemma 08NR: focused-seeds/AGP/sources/stacks-snapshot-20261011/tags.txt, source lines 10999--10999. Navigation pointer into the Stacks Project; the standard fact it supplies is applied inside the counted proof span and the argument resting on it stays there.\par\medskip
\ifdefined\hypertarget\hypertarget{sourcepointer-2}{}\fi
\noindent\textit{Source reference.} Original lemma 08NN: focused-seeds/AGP/sources/stacks-snapshot-20261011/tags.txt, source lines 10996--10996. Navigation pointer into the Stacks Project; the standard fact it supplies is applied inside the counted proof span and the argument resting on it stays there.\par\medskip
\ifdefined\hypertarget\hypertarget{sourcepointer-3}{}\fi
\noindent\textit{Source reference.} Original lemma 08NS: focused-seeds/AGP/sources/stacks-snapshot-20261011/tags.txt, source lines 11000--11000. Navigation pointer into the Stacks Project; the standard fact it supplies is applied inside the counted proof span and the argument resting on it stays there.\par\medskip
\ifdefined\hypertarget\hypertarget{sourcepointer-4}{}\fi
\noindent\textit{Source reference.} Original lemma 08Q9: focused-seeds/AGP/sources/stacks-snapshot-20261011/tags.txt, source lines 11052--11052. Navigation pointer into the Stacks Project; the standard fact it supplies is applied inside the counted proof span and the argument resting on it stays there.\par\medskip
\begin{lemma}
\label{lemma-identify-pi-shriek}
With notation as above let $P_\bullet$ be a simplicial $A$-algebra
endowed with an augmentation $\epsilon : P_\bullet \to B$.
Assume each $P_n$ is a polynomial algebra over $A$ and $\epsilon$
is a trivial Kan fibration on underlying simplicial sets. Then
$$
L\pi_!(\mathcal{F}) = \mathcal{F}(P_\bullet, \epsilon)
$$
in $D(\textit{Ab})$, resp.\ $D(B)$ functorially in $\mathcal{F}$ in
$\textit{Ab}(\mathcal{C})$, resp.\ $\textit{Mod}(\underline{B})$.
\end{lemma}

\begin{proof}
We will use the criterion of Cohomology on Sites, Lemma
\ref{sites-cohomology-lemma-compute-by-cosimplicial-resolution} to prove this.
Given an object $U = (Q, \beta)$ of $\mathcal{C}$ we have to show that
$$
S_\bullet = \Mor_\mathcal{C}((P_\bullet, \epsilon), (Q, \beta))
$$
is homotopy equivalent to a singleton.
Write $Q = A[E]$ for some set $E$ (this is possible by our choice of
the category $\mathcal{C}$). We see that
$$
S_\bullet = \Mor_{\textit{Sets}}((E, \beta|_E), (P_\bullet, \epsilon))
$$
Let $*$ be the constant simplicial set on a singleton. For $b \in B$
let $F_{b, \bullet}$ be the simplicial set defined by the cartesian
diagram
$$
\xymatrix{
F_{b, \bullet} \ar[r] \ar[d] & P_\bullet \ar[d]_\epsilon \\
{*} \ar[r]^b & B
}
$$
With this notation $S_\bullet = \prod_{e \in E} F_{\beta(e), \bullet}$.
Since we assumed $\epsilon$ is a trivial Kan fibration we see that
$F_{b, \bullet} \to *$ is a trivial Kan fibration
(Simplicial, Lemma \ref{simplicial-lemma-trivial-kan-base-change}).
Thus $S_\bullet \to *$ is a trivial Kan fibration
(Simplicial, Lemma \ref{simplicial-lemma-product-trivial-kan}).
Therefore $S_\bullet$ is homotopy equivalent to $*$
(Simplicial, Lemma \ref{simplicial-lemma-trivial-kan-homotopy}).
\end{proof}

\end{document}
